\section{Proofs}

\subsection{Class inclusions and a common K\"ahler background}

Choose a smooth closed representative \(\theta\) of a pseudoeffective
class \(\alpha\), and a K\"ahler form \(\omega\). For sufficiently
large \(M\), \(M\omega-\theta\) is K\"ahler, by compactness of
\(X\). Thus \(\alpha\) maps into the K\"ahler class
\(\kappa=M[\omega]\). A single sufficiently large \(M\) works
for any finite collection of smooth representatives. This also supplies
common upper classes when defining or comparing direct-limit elements.

The DXZ class inclusion is obtained by retaining the later perturbation
components of the inverse system
\cite[Definition~3.5, Lemma~3.7 and (3.5)]{Xia23}. At every strict-interior
parameter its transition preserves every multiplier ideal, and it preserves
the endpoint. Consequently it preserves the literal expression
\begin{equation}\label{eq:p14op-raw}
 F_\Gamma(v)=\sup_{\tau<c_\Gamma}
 \left(\tau-\sup_{m\geq1}m^{-1}
                 v\bigl(\mathcal I(m\Gamma_\tau)\bigr)\right)
\end{equation}
at every semivaluation, not just at divisorial points. Here one uses any
perturbation component, as justified by the fixed P14.2 input. Changes
of a smooth representative use the usual normalized-envelope
identification; addition of a smooth function does not change multiplier
ideals. The raw function therefore depends on the class object, not this
choice of representative.

For comparison, the MP cones increase when a K\"ahler class is added.
On a common test configuration a relatively nef class
\(\kappa_{\mathcal X}+D\) remains relatively nef after adding
the pullback of a K\"ahler class from \(X\). Hence every rational PL
model function positive for \(\kappa\) is positive for the increased
class. The same inclusion holds for their decreasing-net completions.
This proves the inclusions needed in \eqref{eq:p14op-mp-qpsh} directly
from the MP definition \cite[Section~4.1 and Definition~4.3.1]{MP24}.

\begin{proof}[Proof of \cref{thm:p14op-global-realization}]
Embed \(\Gamma\in N(\alpha)\) into \(N(\kappa)\), with both
\(\kappa\) and \(\kappa-\alpha\) K\"ahler. Its raw function
is unchanged. The assumed full K\"ahler comparison makes that same
function MP psh in \(\kappa\). It is nontrivial, since its value at
\(o\) is the finite endpoint. Upper semicontinuity and the upper bound
follow from the assumed P14.2 theorem. If \(q\leq v\) on all coherent
ideals, the inner supremum in \eqref{eq:p14op-raw} increases, so the
whole function decreases. This proves the asserted order reversal.

Compatibility with class inclusions makes the map well-defined on the
DXZ direct limit. Any member of the union in
\eqref{eq:p14op-mp-qpsh} belongs to some K\"ahler MP cone; the full
comparison in that class supplies a DXZ preimage, proving surjectivity.
To compare two DXZ elements, embed them into one common K\"ahler class.
The order isomorphism in that class shows that pointwise comparison, or
equality, of their raw functions is precisely comparison, or equality,
of those representatives. This is the direct-limit order, and proves both
injectivity and order reflection. For a fixed pseudoeffective class,
the same assertion also follows from the source divisorial order criterion
\cite[Lemma~4.2]{Xia23}. No identification of the fixed-class image with
an unchanged MP model completion is used.
\end{proof}

\subsection{An all-point comparison with a larger background}

We recall explicitly the comparison input to avoid a dense-set shortcut.
If \(G:X^\beth\to[-\infty,b]\) is usc and reverses the ideal order,
and if \(\phi\in\PSH_{\mathrm{MP}}(\kappa)\) has the same values
as \(G\) on \(D\), then \(G=\phi\). Indeed, the independently
audited dominated-density input gives a net \(q_\lambda\in D\)
with \(q_\lambda\leq v\) and \(q_\lambda\to v\).
Thus \(G(v)\leq G(q_\lambda)=\phi(q_\lambda)\), and upper
semicontinuity of \(\phi\) gives \(G(v)\leq\phi(v)\), including
when \(\phi(v)=-\infty\). The reverse inequality is MP's comparison
theorem against an arbitrary usc function
\cite[Theorem~4.3.8]{MP24}. Positivity of \(G\) is not assumed.

\begin{proof}[Proof of \cref{thm:p14op-all-point-operations}]
For each operation the source proves existence of its DXZ result and the
displayed identity on \(D\), including the endpoint at \(o\): use
\cite[Proposition~4.4, Corollary~4.8, Definition~4.9,
Lemma~4.10 and Lemma~4.17]{Xia23}. The output classes for positive
scaling and addition are \(a\alpha\) and \(\alpha+\beta\).
These are pseudoeffective. By \cref{thm:p14op-global-realization},
each output raw function is MP psh for some sufficiently large K\"ahler
background class. It remains to check the hypotheses on the right-hand
side of each identity, then apply the preceding comparison.

Finite translations, positive scalar multiples and finite maxima of usc
functions are usc and preserve reversal of the ideal order. They also
preserve the existence of a finite upper bound. For addition, two functions
with finite upper bounds have an usc sum in the extended-real convention:
at a point where both are finite use the usual two upper estimates; at a
point where one is \(-\infty\), use its arbitrarily low local upper
bound and a finite upper bound for the other. Their sum reverses the ideal
order and is bounded above by the sum of the two endpoint bounds. This
proves \eqref{eq:p14op-translation}--\eqref{eq:p14op-maximum}
by the comparison just recalled. No separate positivity claim about these
extended-real expressions was needed in this deduction.

For the decreasing net, fix \(i_0\). Its tail \(i\geq i_0\) is
cofinal, so its infimum is the infimum of the entire net. Every function
in that tail is bounded above by \(c_{\Gamma_{i_0}}\). The pointwise
infimum \(G=\inf_i F_{\Gamma_i}\) is usc, since
\(\{G<t\}=\bigcup_i\{F_{\Gamma_i}<t\}\), and reverses the
ideal order. The finite lower endpoint assumption is precisely what is
needed for existence of the DXZ infimum. The source gives both its
divisorial values and its endpoint as the corresponding infima. Hence
\(G\) agrees with the raw output on \(D\). The all-point comparison
proves \eqref{eq:p14op-infimum}. This argument neither interchanges the
two suprema in \eqref{eq:p14op-raw} with an arbitrary infimum, nor
extracts a sequence from the indexing net.
\end{proof}

\subsection{Boundary checks}

In a K\"ahler class, the source maximum with the constant object
\(-j\) now has the literal raw value
\[
 F_{\Gamma\vee(-j)}(v)=\max\{F_\Gamma(v),-j\}.
\]
At a point where \(F_\Gamma(v)=-\infty\), these values are \(-j\)
and decrease to \(-\infty\). This illustrates why proper-support points
and infinite negative values cannot be discarded. We used a K\"ahler
background for the constant object; the example does not assert that the
constant zero metric belongs to every big non-nef fixed-class space.

The theorem concerns actual values and their elementary operations.
It makes no assertion that an increasing supremum equals its usc
regularization, that continuous-obstacle envelopes are continuous, or that
passing to a larger background preserves a class-specific Monge--Amp\`ere
measure. In particular, the general nontrivially valued envelope problem,
the intrinsic non-nef image, and zero-volume energy remain outside its scope.
